\model{Range-Based For Loops}
  \begin{center}
    \small
    \renewcommand{\arraystretch}{1.2}
    \begin{tabular}{p{2.8in} p{2.8in}}
      \textbf{Traditional For Loop} & \textbf{Range-Based For Loop} \\[4pt]
      \begin{minipage}{2.8in}
        \begin{cpplst}
vector<int> scores =
    {92, 78, 85, 91, 63};

for (int i = 0;
     i < scores.size(); i++) {
    cout << scores.at(i) << " ";
}
cout << endl;
        \end{cpplst}
      \end{minipage}
      &
      \begin{minipage}{2.8in}
        \begin{cpplst}
vector<int> scores =
    {92, 78, 85, 91, 63};

for (int score : scores) {
    cout << score << " ";
}
cout << endl;
        \end{cpplst}
      \end{minipage}
      \\[6pt]
      \multicolumn{2}{c}{%
        \textbf{Output:} \quad {\tt 92 78 85 91 63}
      }
    \end{tabular}
  \end{center}

  {\it\large Refer to Model 1 above as your team develops consensus answers
    to the questions below.}

  \quest{15 min}

  \Q Compare the traditional for loop on the left with the
    range-based for loop on the right.
    \begin{enumerate}
      \itemsep 10pt
      \item What three parts appear in the traditional for loop header
        (before the opening brace), separated by semicolons?
        \begin{answer}[0.5in]
          Initialization (\cpp{int i = 0}), condition
          (\cpp{i < scores.size()}), and increment (\cpp{i++}).
        \end{answer}

      \item In the range-based for loop, what symbol separates the loop
        variable from the container being iterated?
        \hfill\ans[1in]{The colon \cpp{:}}

      \item How many times does the body of each loop execute?
        \hfill\ans[2in]{Five times (once per element)}

      \item What is the value of {\tt score} during the third
        iteration of the range-based for loop?
        \hfill\ans[1in]{85}
    \end{enumerate}

  \vskip -20pt

  \Q The range-based for loop declares a variable that takes on the
    value of each element during each iteration.
    \begin{enumerate}
      \itemsep 10pt
      \item In the range-based for loop above, what is the declared
        type and name of this variable?
        \hfill\ans[2.5in]{Type: \cpp{int}, name: \cpp{score}}

      \item If {\tt scores} were a {\tt vector<string>}, what type should
        the loop variable be declared as?
        \hfill\ans[1.5in]{\cpp{string}}

      \item Write a range-based for loop that prints each element of
        a {\tt vector<double>} named {\tt temps} on its own line.
        \begin{answer}[0.75in]
          \begin{cpplst}
for (double temp : temps) {
    cout << temp << endl;
}
          \end{cpplst}
        \end{answer}
    \end{enumerate}

  \vskip -20pt

  \Q Consider the two loops below. \key\\[-2mm]
    \begin{center}
      \small
      \begin{tabular}{p{2.6in} p{2.6in}}
        \textbf{Loop A} & \textbf{Loop B} \\[4pt]
        \begin{minipage}{2.6in}
          \begin{cpplst}
for (int score : scores) {
    score += 5;
}
          \end{cpplst}
        \end{minipage}
        &
        \begin{minipage}{2.6in}
          \begin{cpplst}
for (int& score : scores) {
    score += 5;
}
          \end{cpplst}
        \end{minipage}
      \end{tabular}
    \end{center}
    \begin{enumerate}
      \itemsep 10pt
      \item What is the only syntactic difference between Loop A and Loop B?
        \hfill\ans[1.5in]{The \cpp{&} after \cpp{int}}

      \item In Loop A, {\tt score} is a copy of each element.  Does
        modifying {\tt score} inside the loop change the original vector?
        \hfill\ans[0.5in]{No}

      \item The \cpp{&} in Loop B declares {\tt score} as a
        {\it reference}, meaning it refers directly to the element
        in the vector.  Does modifying {\tt score} in Loop B change
        the original vector?
        \hfill\ans[0.5in]{Yes}

      \item The file {\tt activity28a.cpp} contains both loop styles.
        Run it and record the output of the last two lines.
        \begin{answer}[0.75in]
          {\tt After copy loop:  92 78 85 91 63}\\
          {\tt After ref loop:   97 83 90 96 68}
        \end{answer}
    \end{enumerate}

  \vskip -20pt

  \Q Write a range-based for loop that computes the sum of all
    elements in a {\tt vector<int>} named {\tt values}.  Use a
    reference when appropriate.
    \begin{answer}[1in]
      References are not needed since we are only reading:
      \begin{cpplst}
int sum = 0;
for (int val : values) {
    sum += val;
}
      \end{cpplst}
    \end{answer}
